\subsection{\texorpdfstring{格罗滕迪克群 \(K_{0}(A)\)}{格罗滕迪克群 K\_\{0\}(A)}}

设 \(\mathrm{A}\) 是一个环。有限生成投射 A-模的类所成之集 \(\mathcal{P}(\mathrm{A})\) 是加性的；我们把格罗滕迪克群 \(\mathrm{K}(\mathcal{P}(\mathrm{A}))\) 记作 \(\mathrm{K}_0(\mathrm{A})\)。

对合成律 \((\mathrm{E},\mathrm{E}^{\prime})\mapsto\mathrm{cl}(\mathrm{E}\oplus\mathrm{E}^{\prime})\)，集 \(\mathcal{P}(\mathrm{A})\) 是一个交换幺半群。此外，投射 A-模的每个正合列

\[
0 \longrightarrow \mathrm{E} ^ {\prime} \longrightarrow \mathrm{E} \longrightarrow \mathrm{E} ^ {\prime \prime} \longrightarrow 0
\]

都分裂（II, §2, No.~2, 第 231 页，命题 4），因而 E 同构于 \(E^{\prime} \oplus E^{\prime\prime}\)。于是，映射 \(E \mapsto [E]\)（从 \(\mathcal{P}(A)\) 到 \(K_{0}(A)\)）定义了从幺半群 \(\mathcal{P}(A)\) 的差群到 \(K_{0}(A)\) 的一个同构（VIII, 第 188 页）。

对每个有限生成投射模 P，存在一个有限生成投射模 \(P'\)，使得 \(P \oplus P'\) 是自由的（II, §2, No.~2, 第 232 页，推论 1）。设 E 与 F 是有限生成投射 A-模；由 I, §2, No.~4, 第 18 页的命题 6，\([E] = [F]\) 在 \(K_{0}(A)\) 中成立当且仅当存在一个有限生成自由 A-模 L，使得 A-模 \(E \oplus L\) 与 \(F \oplus L\) 同构。这时我们说 E 与 F 是稳定同构的；

这并不必然蕴含 E 与 F 同构（VIII, 第 207 页，习题 2 及 VIII, 第 210 页，习题 14）。

当环 A 是交换的时，在加群 \(K_{0}(A)\) 上存在一个交换环结构，其乘法由公式 \([\mathrm{E}]_{\mathcal{P}(\mathrm{A})}[\mathrm{F}]_{\mathcal{P}(\mathrm{A})} = [\mathrm{E} \otimes_{\mathrm{A}} \mathrm{F}]_{\mathcal{P}(\mathrm{A})}\) 刻画（VIII, 第 199 页，命题 10）。

注. —— 设 A 是一个半单环。这时每个 A-模都是半单且投射的（VIII, 第 138 页，命题 4）；因此我们有等式

\[
\mathscr {L} \mathscr {F} (\mathrm{A}) = \mathscr {S S} (\mathrm{A}) = \mathscr {P} (\mathrm{A})
\]

（\(\mathcal{LF}(\mathrm{A})\) 与 \(\mathcal{SS}(\mathrm{A})\) 的定义见 No.~4）。于是由这些格罗滕迪克群的定义，我们有 \(\mathrm{K}_0(\mathrm{A}) = \mathrm{R}(\mathrm{A})\)。

例. —— 若每个有限生成投射 A-模都是自由的，则秩定义了一个从 \(K_{0}(A)\) 到 Z 的同构。当 A 是一个主理想整环（VII, §1, No.~3, 第 15 页，推论 3）或 A 是一个局部环（VIII, 第 36 页，推论 6）时，特别就是这种情形。

